\documentclass[11pt]{article}
\usepackage{fontspec}
\usepackage[french,english]{babel}
\usepackage{amsmath}
\usepackage{amssymb}
\usepackage{amsthm}
\usepackage{mathrsfs}
\usepackage{tikz-cd}
\usepackage{enumitem}

\newcommand{\printedpage}[1]{\marginpar{\footnotesize p.~#1}}
\newcommand{\msp}[1]{\marginpar{\footnotesize [#1]}}
\newcommand{\Loc}{\mathrm{Loc}}

\makeatletter
\renewcommand\@makefnmark{\hbox{\textsuperscript{(\@thefnmark)}}}
\renewcommand\@makefntext[1]{\noindent\textsuperscript{\@thefnmark}\,#1}
\makeatother

\begin{document}

\setcounter{footnote}{2}

\printedpage{4}
% head: A. Grothendieck : Dérivateurs, Ch. XVI, Édition des Universités de Montpellier II et Paris VII
where, as we saw in $(a)$, $i_s^X$ and $i_s^Y$ are in $W$. By hypothesis we have $f_s \in W$, hence $s\backslash f \in W$, q.e.d.

\textit{b}) Loc (1, $2_0$, $3'$) $\Longrightarrow$ Loc (1, 2, 3). Let us first prove Loc (2). Applying Loc ($3'$ \textit{a}) to the projection
\[
\boldsymbol{\Delta}^1 \times X \longrightarrow X \qquad \text{(where $X \in \operatorname{Ob}\mathrm{Cat}$),}
\]
which is split fibered with fibers isomorphic to $\boldsymbol{\Delta}^1$, hence $W$-aspheric by Loc ($2_0$), we see that this projection is in $W$. This allows us to apply the standard homotopy arguments (\textit{cf.} XII, p. 1 \ldots) to conclude that a homotopism is in $W$, hence a contractible category is $W$-aspheric, a fortiori a category with a final object (or with an initial object) is $W$-aspheric, \textit{i.e.} we have Loc (2).

Let us prove Loc (3). Since Loc (1, 2, 3) $\Longleftrightarrow$ Loc (1, $2^\circ$, $3^\circ$), it suffices to prove Loc ($3^\circ$), since we have proved Loc ($2^\circ$). Let
\[
f : X \longrightarrow S \qquad \text{$W$-\textit{co}aspheric,}
\]
let us prove $f \in W$. Let us introduce the \textit{fibered} category over $S$ whose fibers are the $s\backslash X$, namely $\Phi_S(X)$, and consider the canonical inclusion
\[
i_X : X \longrightarrow \Phi_S(X).
\]
\msp{page 7}By Loc ($3'$ \textit{b}) it is in $W$. On the other hand, the fibers $s\backslash X$ of $\Phi_S(X)$ are $W$-coaspheric by hypothesis, hence by Loc ($3'$ \textit{a}) $\Phi_S(X) \longrightarrow S$ is in $W$. Since $i_X \in W$, it follows that $(f : X \longrightarrow S) \in W$, q.e.d.

\textit{b} bis) Loc (1, $2_0$, $3'$ bis) $\Longrightarrow$ Loc (1, 2, 3 bis). We have already proved Loc (1, 2, 3) (since Loc ($3'$ bis) contains Loc ($3'$) as a special case). Let us therefore prove Loc (3 bis) assuming all that already. It will again suffice to prove it in dual form, since
\[
\text{Loc (1, 2, 3 bis)} \iff \text{Loc(1, $2^\circ$, 3 bis$^\circ$),}
\]
\msp{page 8}and Loc ($2^\circ$) is known from Loc (1, 2, 3) $\Longleftrightarrow$ Loc (1, $2^\circ$, $3^\circ$). It thus remains to prove that if
\[
f : X \longrightarrow Y \qquad \text{in } \mathrm{Cat}/S
\]
is such that
\[
s\backslash f : s\backslash X \longrightarrow s\backslash Y
\]
is in $W$ for every $s \in \operatorname{Ob} S$, then $f \in W$. Let us then consider the commutative square
\[
\begin{tikzcd}
X \arrow[r, "i_X"] \arrow[d, "f"'] & \Phi_S(X) \arrow[d, "\Phi_S(f)"] \\
Y \arrow[r, "i_Y"] & \Phi_S(Y),
\end{tikzcd}
\]
where $i_X, i_Y \in W$ by Loc ($3'$ bis \textit{b}). The hypothesis on $f$ means that $\Phi_S(f)$ induces $W$-equivalences on the fibers, hence by Loc ($3'$ bis) $\Phi_S(f) \in W$, hence $f \in W$ by Loc (1), q.e.d.
% foot: Version 18 mars 2019

\printedpage{5}
% head: A. Grothendieck : Dérivateurs, Ch. XVI, Édition des Universités de Montpellier II et Paris VII
\section*{2\quad Axioms Loc (4) and Loc (5) (of direct sums, of connectedness)\protect\footnote{We are dealing here with axioms that one can regard as being of a more or less ``trivial'' nature. In all cases known to me, they are indeed trivially satisfied.}}

\begin{description}[leftmargin=7em, labelwidth=6em]
\item[Loc (4)] (Axiom of \textit{direct sums}). For a family $(f_i)_{i\in I}$ of arrows in Cat, if $f = \coprod f_i$, we have $f \in W$ if and only if $f_i \in W$ for every $i \in I$.
\end{description}

This axiom splits into two, into the ``if'' and the ``only if'', which I am going to formulate in the most economical form:

\textbf{Proposition 2.1.} \textit{The axiom} Loc (4) \textit{of sums is equivalent to the conjunction of the following two axioms:}

\begin{description}[leftmargin=7em, labelwidth=6em]
\item[Loc (4 a)] \msp{page 9}Let $(f_i)_{i\in I}$ be a family of arrows in Cat, if the $f_i \in W$, then $f = \coprod f_i \in W$.
\item[Loc (4 b)] Let $S$ be in Cat, and $S_0 \hookrightarrow S$ a direct summand (subcategory both open and closed). Let $X$ be in Cat$/S$. If $X \longrightarrow S$ is in $W$, then so is $X_0 = X|S_0 \longrightarrow S_0$. (In other words, the inclusion $S_0 \hookrightarrow S$ is a ``$W$-fibration'', \textit{i.e.} is in $\mathrm{Fib}_W$.)
\end{description}

\textbf{Corollary 2.2.} \textit{Consider the variant of} Loc (4 a) \textit{obtained by restricting to finite families.} (\textit{It is equivalent to} Loc (4) \textit{when $W$ is stable under} $\varinjlim$ \textit{that are filtering,} cf. \textit{axiom} Loc (8) \textit{below}). \textit{This variant is equivalent to}

\begin{description}[leftmargin=7.5em, labelwidth=6.5em]
\item[Loc (4 a$_0$)] Every object $X$ of Cat is ``$W$-cofibrant'', \textit{i.e.} $\emptyset \longrightarrow X$ is in $\mathrm{Cof}_W$, \textit{i.e.} for every $f : Y \longrightarrow Y'$ which is in $W$, the corresponding morphism $\mathrm{id}_X \amalg f : X \amalg Y \longrightarrow X \amalg Y'$ is in $W$.
\end{description}

\textbf{Reminder 2.3.} In any case, if $W$ is a basic localizer (which we shall assume from now on, unless otherwise stated), the \textit{finite} analogue of Loc (4 a) for products is true, \textit{i.e.} if we have a \textit{finite} family of arrows $(f_i)_{i\in I}$ in Cat, if the $f_i$ are in $W$, then $f = \prod f_i \in W$. \texttt{[For this, one must assume that $W$ satisfies the localization axiom, \textit{strong} form, Loc 3 bis.]} The analogue for an infinite family is satisfied for practically no basic localizer.

\msp{page 10}\textbf{Proposition 2.4.} \textit{Let $W$ be a basic localizer. Then finite products exist in} $\mathrm{Hot}_W$\textit{, and the functor} $\mathrm{hot}_W : \mathrm{Cat} \longrightarrow \mathrm{Hot}_W \overset{\textit{def}}{=} (\mathrm{Cat})W^{-1}$ \textit{commutes with them. If $W$ satisfies} Loc (4 a) (resp. Loc (4 a$_0$)), \textit{then arbitrary sums} (resp. \textit{finite}) \textit{exist in} $\mathrm{Hot}_W$\textit{, and the functor} $\mathrm{hot}_W$ \textit{commutes with them. Moreover, in this case sums are distributive with respect to products (since this is so in} Cat).

The proof follows simply from the general result:

\textbf{Lemma 2.5.} \textit{Let $(\mathcal{M}, W)$ be a category of models ($W$ any subset of} $\mathrm{Fl}\,\mathcal{M}$). \textit{Suppose that products} (resp. \textit{finite products,} resp. \textit{binary products}) \textit{exist in $\mathcal{M}$, and that moreover for every family $(f_i)_{i\in I}$ of arrows of $\mathcal{M}$} (resp. \textit{every finite family,}
% foot: Version 18 mars 2019
\printedpage{6}
% head: A. Grothendieck : Dérivateurs, Ch. XVI, Édition des Universités de Montpellier II et Paris VII
resp. \textit{binary family}) \textit{$f_i \in W$ for every $i \in I$ implies $f = \prod f_i \in W$. Then products of the type envisaged exist in} $\mathrm{Ho}_W = \mathcal{M}W^{-1}$\textit{, and the functor} $\mathrm{ho}_W : \mathcal{M} \longrightarrow \mathcal{M}W^{-1}$ \textit{commutes with them. Finally, if $\mathcal{M}$ admits a final object $e$, then $e$ is also a final object of} $\mathrm{Ho}_W$.

\texttt{[It is not clear that the assertion is true in this generality for \textit{infinite} products; it is true if $W$ is the class of weak equivalences of a \textit{Quillen model category.}]}

\msp{page 11}We deduce from this the dual statements concerning direct sums and initial objects.

The proof is a special case of the following very useful result, applied to the pair of adjoint functors
\[
\begin{tikzcd}
\mathcal{M} \arrow[r, shift left, "\mathrm{diag}"] & \mathcal{M}^I \arrow[l, shift left, "\text{\scriptsize\shortstack{products\\ in $\mathcal{M}$\\ indexed by $I$}}"]
\end{tikzcd}
\qquad\qquad :
\]

\textbf{Lemma 2.6}\footnote{\textit{cf.} XII, lemma on page 248.}. \textit{Let $(\mathcal{M}, W)$, $(\mathcal{M}', W')$ be two categories of models, and}
\[
\begin{tikzcd}
\mathcal{M} \arrow[r, shift left, "\varphi"] & \mathcal{M}' \arrow[l, shift left, "\psi"]
\end{tikzcd}
\]
\textit{a pair of adjoint functors compatible with the localizers,} i.e. \textit{such that}
\[
\varphi(W) \subset W', \qquad \psi(W') \subset W.
\]
\textit{Then the localized functors}
\[
\begin{tikzcd}
\mathcal{M}W^{-1} \arrow[r, shift left, "\bar{\varphi}"] & \mathcal{M}'W'^{-1} \arrow[l, shift left, "\bar{\psi}"]
\end{tikzcd}
\]
\textit{are also adjoint, the adjunction arrows for $\bar{\varphi}\bar{\psi}$ and $\bar{\psi}\bar{\varphi}$ being deduced from those for $\varphi\psi$ and $\psi\varphi$ by composition with the localization functors $\mathcal{M} \longrightarrow \mathcal{M}W^{-1}$, $\mathcal{M}' \longrightarrow \mathcal{M}'W'^{-1}$.}

\begin{description}[leftmargin=7em, labelwidth=6em]
\item[Loc (5)] \msp{page 12}(Axiom of \textit{connectedness} or of \textit{non-triviality}).

Let $f \in W$, then $\pi_0(f)$ is bijective.
\end{description}

This property of $W$ is in constant use in practice, when we work with $W$ instead of restricting ourselves to constructing $W$'s and playing around with the axioms and their stability under passage to $\underline{\mathrm{Hom}}(I, M)$'s (where here $\mathcal{M} = \mathrm{Cat}$). On the other hand, we can also regard it as an axiom of \textit{non-triviality} for $\mathrm{Hot}_W$. I recall in this connection:

\textbf{Proposition 2.7.} \textit{Let $W$ be a basic localizer. Equivalent conditions:}
\begin{itemize}[leftmargin=3em]
\item[a)] $W$ \textit{satisfies the connectedness axiom} Loc (5).
\item[b)] \textit{The functor} $\pi_0 : \mathrm{Cat} \longrightarrow \mathrm{Ens}$ \textit{factors through a functor}
% foot: Version 18 mars 2019
\printedpage{7}
% head: A. Grothendieck : Dérivateurs, Ch. XVI, Édition des Universités de Montpellier II et Paris VII
\[
\begin{tikzcd}
\mathrm{Cat} \arrow[r, dotted] \arrow[rr, dotted, bend right, "\pi_0"'] & \mathrm{Hot}_W \arrow[r, "\pi_0"] & \mathrm{Ens}.
\end{tikzcd}
\]
\item[c)] \textit{The category} $\mathrm{Hot}^*_W = \mathrm{Hot}_W\backslash\{\emptyset\}$ \textit{is not equivalent to the one-point category.}
\item[c$'$)] \textit{The category} $\mathrm{Hot}_W$ \textit{is equivalent neither to $e$ nor to $\boldsymbol{\Delta}^1$.}\footnote{\textbf{NB} If $W$ is strongly saturated (condition Loc (1 ter), then the conditions opposite (\textit{i.e.} Loc (5)) are equivalent to this: % [sic] "(condition Loc (1 ter)," -> "(condition Loc (1 ter)),", parenthèse ouvrante non refermée
\begin{itemize}[leftmargin=2.5em]
\item[d)] \textit{$W$ is equal neither to Fl}(Cat) \textit{in its entirety, nor to Fl}(Cat$^*$) $\cup$ $\{\mathrm{id}_\emptyset\}$ (i.e. \textit{Fl}(Cat)$\backslash W$ \textit{is neither empty, nor equal to the set of arrows $\emptyset \longrightarrow X$ with $X \neq \emptyset$}).
\end{itemize}
If $W$ satisfies Loc (7 D$_{W_\omega}$), then $W$ satisfies Loc (5) if and only if $W \neq \mathrm{Fl}(\mathrm{Cat})$ (\textit{cf.} XII, th. 2, p. 20).}
\end{itemize}

For the proof, and for other equivalent forms of the connectedness axiom, \msp{page 13}I refer to XV 9.12 (p. 133 bis). For the name ``connectedness axiom'', I can offer the following explanation 2.8 and 2.9.

\textbf{Corollary 2.8.} \textit{If the connectedness axiom} Loc (5) \textit{is satisfied, then the $W$-aspheric objects are} 0\textit{-connected} (\textit{tautology}), \textit{and the converse is true if $W$ is strongly saturated.} (\textit{Or also, if $W$ satisfies the axiom of fibrations} Loc (7) \textit{below, since by} XII, th. 2, p. 20, \textit{in this case,} Loc (5) \textit{is equivalent to $W \neq \mathrm{Fl}(\mathcal{M})$ \ldots}).

For if the connectedness axiom is not satisfied, then for every nonempty $X$ in Cat, $X \longrightarrow e$ is an isomorphism, hence ($W$ being strongly saturated) $X \longrightarrow e$ is in $W$, \textit{i.e.} $X$ is $W$-aspheric, hence there are $W$-aspheric objects that are not 0-connected. Let us also note: % [sic] "non 0-connexe" -> "non 0-connexes", accord avec "des objets"

\textbf{Corollary 2.9.} \textit{Suppose that $W$ satisfies the axiom} Loc (4) \textit{of direct sums} (\textit{or only} Loc (4 b)). \textit{For $W$ to satisfy the connectedness axiom, it is necessary and sufficient that for every arrow $f : X \longrightarrow Y$ with $Y$} 0\textit{-connected, if $f$ is in $W$, then $X$ is} 0\textit{-connected.} (\textit{This is necessary of course without assuming} Loc (4 b) \textit{or anything at all.})

\msp{page 14}\texttt{[The numbers 2.7, 2.8 and 2.9 appear twice.]}

\textbf{2.7.} Let us now consider the composite functor
\[
\mathrm{Ens} \xrightarrow{I \longmapsto I_{\mathrm{Cat}}} \mathrm{Cat} \xrightarrow{\quad \mathrm{hot}_W \quad} \mathrm{Hot}_W, \qquad I \longmapsto I_{\mathrm{Hot}_W} \text{ or } I_W, \tag*{(2.7.1)}
\]
where the first functor is the one that assigns to the set $I$ the discrete category $I_{\mathrm{Cat}}$ having $I$ as its set of objects. (Most often, we shall identify $I_{\mathrm{Cat}}$ with $I$ in the notation, just as for an ordered set $J$, one identifies $J$ and the ``ordered category'' % [sic] "catégorie ordonné" -> "catégorie ordonnée", accord
defined by $J$. Here $I$ can be regarded as equipped with the discrete order.) We have, moreover
\[
I_{\mathrm{Cat}} \simeq I \times e \ \ (\overset{\text{def}}{=} \text{sum of $I$ copies of $e$}), \tag*{(2.7.2)}
\]
so when Loc (4 a) is satisfied, so $\mathrm{hot}_W$ commutes with direct sums, we shall have a canonical isomorphism
\[
\begin{aligned}
I_W \simeq I \times e_W \quad &\text{(where $e_W$ is the final object of $\mathrm{Hot}_W$,} \\
&\text{\textit{i.e.} $e$ regarded as an object of $\mathrm{Hot}_W$).}
\end{aligned} \tag*{(2.7.3)}
\]
% foot: Version 18 mars 2019

\printedpage{8}
% head: A. Grothendieck : Dérivateurs, Ch. XVI, Édition des Universités de Montpellier II et Paris VII
\textbf{Proposition 2.8.} \textit{Suppose the connectedness axiom} Loc (5) \textit{is satisfied, whence a functor}
\[
\pi_0 : \mathrm{Hot}_W \longrightarrow \mathrm{Ens}. \tag*{(2.8.1)}
\]
\textit{This functor is} left adjoint \textit{to the functor $I \longmapsto I_W$} (2.7.1), i.e. \textit{we have a bijection functorial in $(X, I)$ in} $\mathrm{Hot}_W \times \mathrm{Ens}$:
\msp{page 15}
\[
\mathrm{Hom}_{\mathrm{Ens}}(\pi_0(X), I) \simeq \mathrm{Hom}_W(X, I_W). \tag*{(2.8.2)}
\]
\textit{The functor $\pi_0$ is also} right adjoint \textit{to the} same \textit{functor $I \longmapsto I_W$,} i.e. \textit{we have a functorial bijection} % [sic] "un bijection" -> "une bijection", accord
\[
\mathrm{Hom}_{\mathrm{Ens}}(I, \pi_0(X)) \simeq \mathrm{Hom}_W(I_W, X). \tag*{(2.8.3)}
\]

Taking $I = e$ (one-point set), we find the:

\textbf{Corollary 2.9.} \textit{We have the isomorphism functorial in $X$ in} $\mathrm{Hot}_W$
\[
\boxed{\mathrm{Hom}_W(e_W, X) \simeq \pi_0(X).} \tag*{(2.8.4)}
\]

\textbf{Comment 2.10.} The adjunction formula (2.8.2) is, moreover, familiar, whenever we have a category $\mathcal{H}$ (such as Cat, where $\mathrm{Hot}_W$ \textit{etc.}) % [sic] "où" -> "ou", conjonction de coordination
equipped with a functor $\pi_0$, $\pi_0(X)$ corresponding to the intuition of a decomposition of an object into ``connected components''. (I am not going to spell out here the general conditions under which one defines such a functor.) On the other hand, the adjunction in the reverse direction of these same two functors is a phenomenon quite special to categories \msp{page 16}of the type $\mathrm{Hot}_W$. If we take into account the description of $I_W$ as $I \times e_W$, the unusual adjunction formula (2.8.3) is equivalent, of course, to formula (2.8.4) - it is deduced from it by ``raising to the power $I$''. We can reverse the steam, and in a general category $\mathcal{H}$, \textit{define} a functor $\pi_0 : \mathcal{H} \longrightarrow \mathrm{Ens}$ by formula (2.8.4), % [sic] "par le formule" -> "par la formule", accord
\[
\pi_0(X) \overset{\text{def}}{=} \mathrm{Hom}(e_{\mathcal{H}}, X) \qquad (= \Gamma_{\mathcal{H}}(X), \text{ with a familiar notation}),
\]
where $e_{\mathcal{H}}$ is a final object of $\mathcal{H}$. We shall therefore have, putting $I_{\mathcal{H}} = I \times e_{\mathcal{H}}$, the formula
\[
\mathrm{Hom}(I_{\mathcal{H}}, X) \xrightarrow{\ \sim\ } \mathrm{Hom}(I, \pi_0(X))
\]
(\textit{i.e.} formula (2.8.3)), saying that we have a pair $(I_{\mathcal{H}}, \pi_0)$ of adjoint functors. Suppose on the other hand that for every set $I$, the canonical adjunction map
\[
I \longrightarrow \pi_0(I_{\mathcal{H}}) = \mathrm{Hom}_{\mathcal{H}}(e_{\mathcal{H}}, e_{\mathcal{H}} \times I)
\]
is \textit{bijective}\footnote{that tends to be ultrafalse if $\mathcal{H}$ is additive. On the other hand, true if $\mathcal{H}$ is a category with sums distributive with respect to products (or even universal), whose final object is \textit{connected} \ldots}, \textit{i.e.} that $I \longmapsto I_{\mathcal{H}}$ is \textit{fully faithful.} We then define, for every $I$, a canonical map
% foot: Version 18 mars 2019

\end{document}
